\section{Discussão}
\label{sec:discussao}

Os resultados mostram que atualizar apenas as estatísticas do Batch
Normalization não melhorou o desempenho neste experimento. Uma hipótese
para a queda observada é que a normalização pelas estatísticas dos lotes
altere as representações do extrator sem permitir a adaptação de seus
parâmetros de escala e deslocamento. O ganho obtido com
BN treinável sugere que ajustar esses parâmetros contribuiu
para adaptar o modelo ao Food-101, mesmo com as convoluções congeladas.
Entretanto, os experimentos realizados não permitem determinar isoladamente
a causa da queda em BN estatístico.

As curvas também mostram diferenças na evolução do aprendizado. No
modelo de referência, a perda de validação atinge seu mínimo na quinta época;
em BN estatístico, na quarta. Em ambos, a perda de validação termina
acima do mínimo observado. Em
BN treinável e BN treinável + SD, a menor perda
de validação ocorre na décima
época, indicando que o treinamento ainda apresentou melhora nesse critério
até o final do período avaliado.

O Stochastic Depth acrescentou um ganho de 0,50 ponto percentual em
acurácia top-1 de teste, sem aumentar a quantidade de parâmetros treináveis.
Sua perda final de validação também foi menor que a de
BN treinável (0,9866 contra 1,0045).

A comparação está limitada a uma execução por configuração, com semente
fixa e duração de 10 épocas. Assim, não é possível estimar a variabilidade
dos resultados entre execuções nem afirmar que o ganho adicional do
Stochastic Depth seja consistente entre sementes. As conclusões se
restringem às configurações e às condições avaliadas.
